\documentclass[11pt]{article}
\usepackage[margin=1in]{geometry}
\usepackage{amsmath,amssymb}
\usepackage{graphicx}
\usepackage{hyperref}
\usepackage{listings}
\usepackage{xcolor}

\lstset{
  basicstyle=\ttfamily\small,
  backgroundcolor=\color{gray!10},
  frame=single,
  breaklines=true
}

\title{Radiative Corrections with \texttt{aao\_rad} and \texttt{aao\_norad}:\\
An Example of the Calculation as a Function of $\phi$}
\author{V.~Kubarovsky}
\date{\today}

\begin{document}
\maketitle

\begin{abstract}
This note describes how to use the two versions of the \texttt{aao\_gen}
event generator, \texttt{aao\_rad} (with internal radiative effects) and
\texttt{aao\_norad} (Born-level, non-radiative), to calculate the radiative
correction factor for exclusive single-pion electroproduction. As a working
example, the correction is evaluated as a function of the azimuthal angle
$\phi$ of the hadronic plane. The same procedure applies to any kinematic
variable ($Q^2$, $W$, $t$, $\cos\theta^*$, \dots).
\end{abstract}

\section{The generators}

The repository \texttt{aao\_gen} contains two Fortran generators:
\begin{itemize}
  \item \texttt{aao\_norad} --- generates events according to the Born
        (non-radiative) cross section of the chosen physics model
        (A0, MAID98, MAID2000);
  \item \texttt{aao\_rad} --- generates events including internal radiative
        effects (vertex corrections, vacuum polarization, and real photon
        emission) following the Mo--Tsai formalism applied to the same
        underlying model.
\end{itemize}

Both generators use accept--reject sampling, so the output events are
\emph{unweighted}. At the end of the run each generator prints the
\textbf{integrated cross section} over the generation region,
\begin{lstlisting}
 Integrated cross section = sig_int  sig_sum
\end{lstlisting}
computed by Monte Carlo integration as
\begin{equation}
  \sigma_{\rm tot} \;=\;
  \frac{\sum_k \sigma_k}{N_{\rm tries}}\,(4\pi)^2\,
  \bigl(2\pi\,\Delta Q^2\,\Delta E'\bigr),
\end{equation}
i.e.\ the total cross section \emph{within the kinematic limits specified in
the input file} ($Q^2$ range, scattered-electron energy range, $W$ range).

\section{From event samples to differential cross sections}

Since the events are unweighted, a sample of $N$ events together with the
integrated cross section $\sigma_{\rm tot}$ fully defines the differential
cross section in any variable. For a histogram in $\phi$ with bin width
$\Delta\phi$ and $n_i$ events in bin $i$,
\begin{equation}
  \frac{d\sigma}{d\phi}\bigg|_i \;=\;
  \sigma_{\rm tot}\,\frac{n_i}{N\,\Delta\phi}.
  \label{eq:dsig}
\end{equation}

\section{The radiative correction factor}

The radiative correction factor in bin $i$ is defined as the ratio of the
radiated to the Born differential cross sections,
\begin{equation}
  RC(\phi_i) \;=\;
  \frac{d\sigma_{\rm rad}/d\phi}{d\sigma_{\rm norad}/d\phi}
  \;=\;
  \frac{\sigma_{\rm rad}}{\sigma_{\rm norad}}\cdot
  \frac{n_i^{\rm rad}/N_{\rm rad}}{n_i^{\rm norad}/N_{\rm norad}}.
  \label{eq:rc}
\end{equation}
For equal statistics, $N_{\rm rad}=N_{\rm norad}$ (e.g.\ $100\,000$ events
each), this reduces to the simple prescription:
\begin{enumerate}
  \item fill the \texttt{aao\_norad} histogram $h_{\rm norad}(\phi)$ with
        weight $1$;
  \item fill the \texttt{aao\_rad} histogram $h_{\rm rad}(\phi)$ with the
        constant weight
        \begin{equation}
          w \;=\; \frac{\sigma_{\rm rad}}{\sigma_{\rm norad}},
        \end{equation}
        where $\sigma_{\rm rad}$ and $\sigma_{\rm norad}$ are the two
        integrated cross sections printed at the end of the runs;
  \item divide bin by bin:
        \begin{equation}
          RC(\phi) \;=\; \frac{h_{\rm rad}(\phi)}{h_{\rm norad}(\phi)}.
        \end{equation}
\end{enumerate}

The overall normalization $\sigma_{\rm rad}/\sigma_{\rm norad}$ carries the
bulk of the correction (vertex correction, vacuum polarization, soft-photon
emission), while the bin-by-bin shape ratio adds the kinematic migration
caused by hard-photon emission.

\section{The internal missing-mass cut in \texttt{aao\_rad} (\texttt{mm\_cut})}
\label{sec:mmcut}

The \texttt{aao\_rad} input file contains the parameter \texttt{mm\_cut}
(the line labeled ``a limit on the error in (mm)**2''), which is applied
\emph{inside the generator, during event generation}:
\begin{lstlisting}
if (abs(mm2 - mm_exp) .gt. mm_cut) go to 20   ! reject, regenerate
sig_tot = sig_tot + sigr                      ! only surviving events count
\end{lstlisting}

Here \texttt{mm2} is the \emph{experimental} missing mass squared, computed
the way an experimentalist computes it --- from the nominal beam energy, the
scattered electron, and the detected hadron (the proton for $\pi^0$/$\eta$,
the pion for $\pi^+$), \emph{ignoring the radiated photon}. For the $\pi^0$
channel it is $MM^2(ep\to e'pX)$. The reference value \texttt{mm\_exp} is
the nominal missing mass squared of the channel: $m^2_{\pi^0}$ for
\texttt{epirea=1}, $m^2_p\simeq m^2_n$ for \texttt{epirea=3}, $m^2_\eta$
for \texttt{epirea=5}. When a hard photon is emitted, \texttt{mm2} shifts
away from \texttt{mm\_exp}; this shift \emph{is} the radiative tail.

Since the rejection happens \emph{before} the cross section is accumulated,
\textbf{the integrated cross section $\sigma_{\rm rad}$ printed at the end
of the run is the cross section within the window
$|MM^2 - MM^2_{\rm exp}| < \texttt{mm\_cut}$}. The parameter therefore
plays exactly the role of the inelasticity cut $v_{\rm cut}$ in other
radiative-correction codes: it defines how much of the radiative tail the
sample --- and $\sigma_{\rm rad}$ --- contains, and hence the size of the
correction itself.

Practical rules:
\begin{enumerate}
  \item The radiative correction is only defined \emph{with respect to a
        specific missing-mass cut}; \texttt{mm\_cut} must be consistent with
        the $MM^2$ cut applied to the data.
  \item Recommended: set \texttt{mm\_cut} \emph{looser} than the analysis
        cut (the data $MM^2$ is smeared by detector resolution, so events
        migrate across the cut boundary in both directions), and apply the
        exact analysis cut offline on the \emph{reconstructed} $MM^2$ after
        GEMC, with the same code as for the data. The normalization stays
        correct because
        $d\sigma = \sigma_{\rm printed}\times n_{\rm pass}/N_{\rm generated}$.
        Do not set it absurdly loose, or most of the generated events fall
        in the tail that the analysis cut later removes.
  \item \textbf{Never} set \texttt{mm\_cut} tighter than the data cut: the
        part of the radiative tail present in the data would then be missing
        from the Monte Carlo, and the correction would be underestimated
        with no way to recover it offline.
  \item \texttt{aao\_norad} has no radiated photon, so at generator level
        its $MM^2$ is exact and no such parameter exists there. After
        detector smearing, apply the same reconstructed-$MM^2$ cut to both
        samples.
\end{enumerate}

\section{Applying the correction to data}

The experimental (radiated) cross section is corrected to the Born level by
dividing by the $RC$ factor evaluated in the same bin with the same cuts:
\begin{equation}
  \frac{d\sigma_{\rm Born}}{d\phi}\bigg|_{\rm data}
  \;=\;
  \frac{1}{RC(\phi)}\,
  \frac{d\sigma_{\rm exp}}{d\phi}\bigg|_{\rm data}.
\end{equation}

\section{Practical remarks}

\begin{itemize}
  \item \textbf{Identical generation limits.}
        $\sigma_{\rm tot}$ is integrated over the generation region, so both
        runs must use exactly the same input limits ($Q^2$, $E'$, $W$
        ranges); otherwise the ratio $\sigma_{\rm rad}/\sigma_{\rm norad}$ is
        meaningless.
  \item \textbf{Reconstruct kinematics as in the data.}
        For radiated events, $\phi$, $Q^2$, $W$ must be computed from the
        scattered electron (and detected hadrons), not from the true
        hadronic-vertex values. The difference between the two is precisely
        the radiative migration being corrected.
  \item \textbf{The exclusivity cut defines the correction.}
        The size of $RC$ depends on how much of the radiative tail the
        missing-mass (or $W$) cut accepts. The cut applied to the Monte
        Carlo must match the cut applied to the data exactly; see
        Sec.~\ref{sec:mmcut} for the internal \texttt{mm\_cut} parameter of
        \texttt{aao\_rad}.
  \item \textbf{Generator level vs.\ after detector simulation.}
        Taking the ratio after GEMC is legitimate: the acceptance largely
        cancels in the ratio, and the surviving radiation-induced acceptance
        change is genuinely part of the correction. The cleaner standard
        split is to take $RC$ from the generator-level ratio with
        data-equivalent cuts and the acceptance from the radiated sample
        alone. Either way, do not correct for the acceptance twice.
  \item \textbf{Bin in the other variables.}
        $RC$ depends on $Q^2$ and $W$, so an inclusive $RC(\phi)$ is an
        acceptance-weighted average; evaluate it within the $(Q^2,W)$ or
        $(Q^2,t)$ bins of the analysis.
  \item \textbf{Same polarization flag.}
        \texttt{aao\_rad} modulates the cross section with polarization and
        interference terms; use the same \texttt{flag\_ehel} in both runs.
  \item \textbf{Statistics.}
        With $100\,000$ events and 20 $\phi$ bins one has $\sim 5000$
        events per bin at generator level, i.e.\ $\sim 2\%$ statistical
        uncertainty on $RC$ per bin. After detector acceptance the available
        statistics may be substantially smaller; generate more events if the
        ratio is formed after GEMC.
\end{itemize}

\section{Summary}

\begin{equation}
  \boxed{\;
  RC(\phi) \;=\;
  \frac{\sigma_{\rm rad}}{\sigma_{\rm norad}}\cdot
  \frac{h_{\rm rad}(\phi)}{h_{\rm norad}(\phi)}
  \qquad\Longrightarrow\qquad
  \sigma_{\rm Born} \;=\; \frac{\sigma_{\rm exp}}{RC}
  \;}
\end{equation}
where the histograms are filled with equal numbers of unweighted events
under identical generation limits and identical analysis cuts.

\end{document}
